\documentclass[a4paper]{arrowhead}

\usepackage[yyyymmdd]{datetime}
\usepackage{etoolbox}
\usepackage[utf8]{inputenc}
\usepackage[T1]{fontenc}
\usepackage{multirow}
\usepackage{longtable}
\usepackage{tikz}
\usepackage{adjustbox}
\AtBeginEnvironment{longtable}{\small}
\BeforeBeginEnvironment{tikzpicture}{\begin{adjustbox}{max width=\textwidth}}
\AfterEndEnvironment{tikzpicture}{\end{adjustbox}}
\usetikzlibrary{positioning,shapes.multipart,arrows.meta}

\renewcommand{\dateseparator}{-}
\setlength{\emergencystretch}{3em}
\lstset{language=HTTP,breakatwhitespace=false,columns=fullflexible}

%% Special references
\newcommand{\mref}[2][]{{\textcolor{ArrowheadPurple}{\hyperref[sec:model:#2]{\ifstrempty{#1}{#2}{#1}}}}}
\newcommand{\oref}[1]{{\textcolor{ArrowheadBlue}{\hyperref[sec:operations:#1]{#1}}}}
\newcommand{\pdef}[1]{{\textcolor{ArrowheadGrey}{#1\label{sec:model:primitives:#1}\label{sec:model:primitives:#1s}\label{sec:model:primitives:#1es}}}}
\newcommand{\pref}[1]{{\textcolor{ArrowheadGrey}{\hyperref[sec:model:primitives:#1]{#1}}}}

% as template_sd.tex; the input and output arguments are passed already formatted (\mref or \pref)
\newrobustcmd\osubsection[3]{
  \addtocounter{subsection}{1}
  \addcontentsline{toc}{subsection}{\protect\numberline{\thesubsection}operation \textcolor{ArrowheadBlue}{#1}}
  \subsection*{
    \thesubsection\quad
    operation
    \textcolor{ArrowheadBlue}{#1}
    (\notblank{#2}{#2}{})
    \notblank{#3}{: #3}{}
  }
  \label{sec:operations:#1}
}
\newrobustcmd\msubsection[2]{
  \addtocounter{subsection}{1}
  \addcontentsline{toc}{subsection}{\protect\numberline{\thesubsection}#1 \textcolor{ArrowheadPurple}{#2}}
  \subsection*{\thesubsection\quad#1 \textcolor{ArrowheadPurple}{#2}}
  \label{sec:model:#2} \label{sec:model:#2es}
}
%%

\begin{document}
\sloppy

%% Arrowhead Document Properties
\ArrowheadTitle{serviceOrchestrationSimpleStoreManagement}
\ArrowheadServiceID{serviceOrchestrationSimpleStoreManagement}
\ArrowheadType{Service Description}
\ArrowheadTypeShort{SD}
\ArrowheadVersion{5.2.2}
\ArrowheadDate{2026-10-22}
\ArrowheadAuthor{Jerker Delsing}
\ArrowheadStatus{RELEASE CANDIDATE}
\ArrowheadContact{jerker.delsing@ltu.se}
\ArrowheadFooter{\href{www.arrowhead.eu}{www.arrowhead.eu}}
\ArrowheadSetup
%%

%% Front Page
\begin{center}
  \vspace*{1cm}
  \huge{\arrowtitle}

  \vspace*{0.2cm}
  \LARGE{\arrowtype}
  \vspace*{1cm}

  \Large{Service ID: \textit{"\arrowid"}}

  \vspace*{\fill}

  % Front Page Image

  \vspace*{1cm}
  \vspace*{\fill}

  % Front Page Abstract
  \begin{abstract}
    This document describes the \texttt{service\allowbreak{}Orchestration\allowbreak{}Simple\allowbreak{}Store\allowbreak{}Management} service, version 1.0.0, produced by the SimpleStoreServiceOrchestration core system of Arrowhead 5.2.2.
  \end{abstract}

  \vspace*{1cm}
\end{center}

\newpage
%%

%% Table of Contents
\tableofcontents
\newpage
%%

\section{Overview}
\label{sec:overview}

This document describes the \texttt{service\allowbreak{}Orchestration\allowbreak{}Simple\allowbreak{}Store\allowbreak{}Management} service, which provides operator management of the orchestration store: the static rules telling which service instance a consumer shall use, with priorities.
It is produced by the SimpleStoreServiceOrchestration system \cite{ssso-sysd} in service version 1.0.0.

The rest of this document is organized as follows.
In Section \ref{sec:operations}, we describe the abstract message operations provided by the service.
In Section \ref{sec:model}, we end the document by presenting the data types used by the mentioned operations.
The concrete interface is specified in the HTTP/TLS/JSON IDD \cite{ssso-serviceOrchestrationSimpleStoreManagement-http-idd}. The MQTT binding is specified in the MQTT/TLS/JSON IDD \cite{ssso-serviceOrchestrationSimpleStoreManagement-mqtt-idd}.

\subsection{Significant Prior Art}

Successor of the Arrowhead 4.x Orchestrator store management.

\subsection{How This Service Is Meant to Be Used}

Operators and support systems such as the Plant\allowbreak{}Description\allowbreak{}Engine create rules (consumer, service instance, priority), query, reprioritize and remove them. \texttt{pull} then returns the rules of the consumer in priority order.

\subsection{Important Delimitations and Dependencies}
\label{sec:delimitations}

Rules are not checked against the ServiceRegistry or Consumer\allowbreak{}Authorization at orchestration time.

\newpage

\section{Service Interface}
\label{sec:operations}

This section describes the interfaces to the \texttt{service\allowbreak{}Orchestration\allowbreak{}Simple\allowbreak{}Store\allowbreak{}Management} service.
In particular, each subsection names an abstract operation, an input type and an output type, in that order.
The input type is named inside parentheses, while the output type is preceded by a colon.
Input and output types are only denoted when accepted or returned, respectively, by the interface in question.

All abstract data types named in this section are defined in Section \ref{sec:model}.

\begin{figure}[ht!]
  \centering
  \begin{tikzpicture}[font=\small]
    \node[draw=ArrowheadBlue, very thick, rectangle split, rectangle split parts=2, align=center,
          rectangle split part fill={ArrowheadBlue!12,white}, text width=9cm, inner sep=6pt] (b)
      {\guillemotleft{}Service\guillemotright{}\\\textbf{serviceOrchestrationSimpleStoreManagement}
        \nodepart{second} \texttt{query}\\ \texttt{remove}\\ \texttt{create}\\ \texttt{modify-\allowbreak{}priorities}};
  \end{tikzpicture}
  \caption{SysML-style block of the \texttt{service\allowbreak{}Orchestration\allowbreak{}Simple\allowbreak{}Store\allowbreak{}Management} service and its operations.}
  \label{fig:SD}
\end{figure}

The following interface operations are available.

\osubsection{query}{\mref[Orchestration\allowbreak{}Simple\allowbreak{}Store\allowbreak{}Query\allowbreak{}Request]{OrchestrationSimpleStoreQueryRequest}}{\mref[Orchestration\allowbreak{}Simple\allowbreak{}Store\allowbreak{}List\allowbreak{}Response]{OrchestrationSimpleStoreListResponse}}

Returns a page of store entries matching the filter.

\osubsection{remove}{\pref{List}$<$\allowbreak{}\pref{String}$>$}{}

Removes store entries by identifier.

\osubsection{create}{\mref[Orchestration\allowbreak{}Simple\allowbreak{}Store\allowbreak{}List\allowbreak{}Request]{OrchestrationSimpleStoreListRequest}}{\mref[Orchestration\allowbreak{}Simple\allowbreak{}Store\allowbreak{}List\allowbreak{}Response]{OrchestrationSimpleStoreListResponse}}

Creates store entries; the service definition is taken from the instance identifier.

\osubsection{modify-priorities}{\mref{PriorityRequest}}{\mref[Orchestration\allowbreak{}Simple\allowbreak{}Store\allowbreak{}List\allowbreak{}Response]{OrchestrationSimpleStoreListResponse}}

Changes the priorities of store entries.

\newpage

\section{Information Model}
\label{sec:model}

Here, all data objects that can be part of the \texttt{service\allowbreak{}Orchestration\allowbreak{}Simple\allowbreak{}Store\allowbreak{}Management} service are listed in alphabetic order.
Note that each subsection, which describes one type of object, begins with the \textit{struct} keyword, which is used to denote a collection of named fields, each with its own data type.
As a complement to the explicitly defined types in this section, there is also a list of implicit primitive types in Section \ref{sec:model:primitives}, which are used to represent things like hashes and identifiers.

\msubsection{struct}{OrchestrationSimpleStoreListRequest}

Store entries to create.


\noindent\begin{tabularx}{\textwidth}{| >{\raggedright\arraybackslash}p{3.9cm} | >{\raggedright\arraybackslash\small}p{5.3cm} | X |} \hline
\rowcolor{gray!33} Field & Type & Description \\ \hline
candidates & \pref{List}$<$\allowbreak{}\mref[Orchestration\allowbreak{}Simple\allowbreak{}Store\allowbreak{}Request]{OrchestrationSimpleStoreRequest}$>$ & Store entries to create. \\ \hline
\end{tabularx}



\msubsection{struct}{OrchestrationSimpleStoreListResponse}

List of store entries.


\noindent\begin{tabularx}{\textwidth}{| >{\raggedright\arraybackslash}p{3.9cm} | >{\raggedright\arraybackslash\small}p{5.3cm} | X |} \hline
\rowcolor{gray!33} Field & Type & Description \\ \hline
entries & \pref{List}$<$\allowbreak{}\mref[Orchestration\allowbreak{}Simple\allowbreak{}Store\allowbreak{}Response]{OrchestrationSimpleStoreResponse}$>$ & Matching records of the current page. \\ \hline
count & \pref{Number} & Total number of matching records (all pages). \\ \hline
\end{tabularx}



\msubsection{struct}{OrchestrationSimpleStoreQueryRequest}

Paged filter for store entries.


\noindent\begin{tabularx}{\textwidth}{| >{\raggedright\arraybackslash}p{3.9cm} | >{\raggedright\arraybackslash\small}p{5.3cm} | X |} \hline
\rowcolor{gray!33} Field & Type & Description \\ \hline
pagination & \mref{Page} & Paging and sorting of the result. \\ \hline
ids & \pref{List}$<$\allowbreak{}\pref{Identifier}$>$ & Store entry identifiers (UUID) to match. \\ \hline
consumer\allowbreak{}Names & \pref{List}$<$\allowbreak{}\pref{String}$>$ & Consumer system names to match. \\ \hline
service\allowbreak{}Definitions & \pref{List}$<$\allowbreak{}\pref{Name}$>$ & Service definition names to match. \\ \hline
service\allowbreak{}Instance\allowbreak{}Ids & \pref{List}$<$\allowbreak{}\pref{Identifier}$>$ & Service instance identifiers to match. \\ \hline
min\allowbreak{}Priority & \pref{Number} & Lower bound of the priority. \\ \hline
max\allowbreak{}Priority & \pref{Number} & Upper bound of the priority. \\ \hline
created\allowbreak{}By & \pref{Name} & Creator system to match. \\ \hline
\end{tabularx}



\msubsection{struct}{OrchestrationSimpleStoreRequest}

One orchestration rule: consumer, service instance and priority.


\noindent\begin{tabularx}{\textwidth}{| >{\raggedright\arraybackslash}p{3.9cm} | >{\raggedright\arraybackslash\small}p{5.3cm} | X |} \hline
\rowcolor{gray!33} Field & Type & Description \\ \hline
consumer & \pref{Name} & Consumer system the rule applies to. \\ \hline
service\allowbreak{}Instance\allowbreak{}Id & \pref{Identifier} & Service instance offered to the consumer. \\ \hline
priority & \pref{Number} & Priority of the rule; 1 is the highest. \\ \hline
\end{tabularx}



\msubsection{struct}{OrchestrationSimpleStoreResponse}

A stored orchestration rule.


\noindent\begin{tabularx}{\textwidth}{| >{\raggedright\arraybackslash}p{3.9cm} | >{\raggedright\arraybackslash\small}p{5.3cm} | X |} \hline
\rowcolor{gray!33} Field & Type & Description \\ \hline
id & \pref{Identifier} & Identifier (UUID) of the store entry. \\ \hline
consumer & \pref{Name} & Name of the consumer system. \\ \hline
service\allowbreak{}Definition & \pref{Name} & Service definition name. \\ \hline
service\allowbreak{}Instance\allowbreak{}Id & \pref{Identifier} & Service instance identifier \texttt{\textless{}System\textgreater{}\textbar{}\allowbreak{}\textless{}service\allowbreak{}Definition\textgreater{}\textbar{}\allowbreak{}\textless{}version\textgreater{}}. \\ \hline
priority & \pref{Number} & Priority of the entry; 1 is the highest. \\ \hline
created\allowbreak{}By & \pref{Name} & Name of the system that created the record. \\ \hline
updated\allowbreak{}By & \pref{Name} & Name of the system that changed the record last. \\ \hline
created\allowbreak{}At & \pref{DateTime} & UTC time of creation. \\ \hline
updated\allowbreak{}At & \pref{DateTime} & UTC time of the last change. \\ \hline
\end{tabularx}



\msubsection{struct}{Page}

Paging and sorting; all fields optional.


\noindent\begin{tabularx}{\textwidth}{| >{\raggedright\arraybackslash}p{3.9cm} | >{\raggedright\arraybackslash\small}p{5.3cm} | X |} \hline
\rowcolor{gray!33} Field & Type & Description \\ \hline
page & \pref{Number} & Zero-based page number. \\ \hline
size & \pref{Number} & Number of records per page (at most \texttt{max.\allowbreak{}page.\allowbreak{}size}, default 1000). \\ \hline
direction & \pref{String} & Sort direction, \texttt{ASC} or \texttt{DESC}. \\ \hline
sort\allowbreak{}Field & \pref{String} & Field to sort by. \\ \hline
\end{tabularx}



\msubsection{struct}{PriorityRequest}

A JSON object mapping store entry identifiers (UUID) to new priorities, e.g. \texttt{\{"123e4567-\allowbreak{}e89b-\allowbreak{}12d3-\allowbreak{}a456-\allowbreak{}426614174000": 5\}}.



\subsection{Primitives}
\label{sec:model:primitives}

Types and structures mentioned throughout this document that are assumed to be available to implementations of this service.
The concrete interpretations of each of these types and structures must be provided by any IDD document claiming to implement this service.

\noindent\begin{tabularx}{\textwidth}{| p{4.6cm} | X |} \hline
\rowcolor{gray!33} Type & Description \\ \hline
\pdef{DateTime} & Pinpoints a specific moment in time (UTC). \\ \hline
\pdef{Identifier} & A string that identifies a record, e.g. a UUID or a composite identifier. \\ \hline
\pdef{List} & An \textit{array} of a known number of items, each having type A. \\ \hline
\pdef{Name} & A string identifier that is intended to be both human and machine-readable. \\ \hline
\pdef{Number} & An integer number. \\ \hline
\pdef{String} & A chain of characters. \\ \hline
\end{tabularx}

\newpage

\bibliographystyle{IEEEtran}
\bibliography{bibliography}

\newpage

\section{Revision History}
\subsection{Amendments}

\noindent\begin{tabularx}{\textwidth}{| p{1cm} | p{3cm} | p{2cm} | X | p{4cm} |} \hline
\rowcolor{gray!33} No. & Date & Version & Subject of Amendments & Author \\ \hline

1 & 2026-10-22 & \arrowversion & Document created from the Arrowhead LaTeX templates for the 5.2.2 Java Spring implementation, including the local extensions of this workspace. & Jerker Delsing \\ \hline
\end{tabularx}

\subsection{Quality Assurance}

\noindent\begin{tabularx}{\textwidth}{| p{1cm} | p{3cm} | p{2cm} | X |} \hline
\rowcolor{gray!33} No. & Date & Version & Approved by \\ \hline

1 & & \arrowversion  &  \\ \hline

\end{tabularx}

\end{document}
